% ============================================================================
%  preamble.tex — Configuration du document
% ============================================================================

% ---------------------------------------------------------------------------
%  Langue et encodage
% ---------------------------------------------------------------------------
\usepackage[T1]{fontenc}
\usepackage[utf8]{inputenc}
\usepackage{lmodern}
\usepackage[french]{babel}
\usepackage[autostyle,french=guillemets]{csquotes}
\frenchsetup{SmallCapsFigTabCaptions=false}

% ---------------------------------------------------------------------------
%  Mise en page
% ---------------------------------------------------------------------------
\usepackage[a4paper,inner=3cm,outer=2.5cm,top=2.5cm,bottom=2.5cm]{geometry}
\usepackage{setspace}
\onehalfspacing
\usepackage{emptypage}
\usepackage{fancyhdr}
\pagestyle{fancy}
\fancyhf{}
\fancyhead[LE,RO]{\thepage}
\fancyhead[RE]{\nouppercase{\leftmark}}
\fancyhead[LO]{\nouppercase{\rightmark}}
\renewcommand{\headrulewidth}{0.4pt}
\fancypagestyle{plain}{%
  \fancyhf{}%
  \fancyfoot[C]{\thepage}%
  \renewcommand{\headrulewidth}{0pt}%
}

% ---------------------------------------------------------------------------
%  Graphiques, tableaux, flottants
% ---------------------------------------------------------------------------
\usepackage{graphicx}
\graphicspath{{figures/}}
\usepackage{amssymb}   % \checkmark
\usepackage{booktabs}
\usepackage{tabularx}
\usepackage{longtable}
\usepackage{multirow}
\usepackage{float}
\usepackage[font=small,labelfont=bf]{caption}
\usepackage{subcaption}
\usepackage{xcolor}
\usepackage{enumitem}

% ---------------------------------------------------------------------------
%  Listings (extraits YAML, Bash, Dockerfile, HCL)
% ---------------------------------------------------------------------------
\usepackage{listings}

\definecolor{codebg}{RGB}{248,248,248}
\definecolor{codeframe}{RGB}{210,210,210}
\definecolor{codekw}{RGB}{0,86,160}
\definecolor{codestr}{RGB}{163,21,21}
\definecolor{codecmt}{RGB}{96,139,78}

\lstdefinestyle{sestyle}{
  backgroundcolor=\color{codebg},
  basicstyle=\ttfamily\footnotesize,
  keywordstyle=\color{codekw}\bfseries,
  stringstyle=\color{codestr},
  commentstyle=\color{codecmt}\itshape,
  numbers=left,
  numberstyle=\tiny\color{gray},
  numbersep=8pt,
  frame=single,
  rulecolor=\color{codeframe},
  breaklines=true,
  breakatwhitespace=true,
  showstringspaces=false,
  tabsize=2,
  captionpos=b,
  literate=%
    {é}{{\'e}}1 {è}{{\`e}}1 {ê}{{\^e}}1 {à}{{\`a}}1
    {ù}{{\`u}}1 {ç}{{\c c}}1 {ô}{{\^o}}1 {î}{{\^\i}}1
}
\lstset{style=sestyle}

% Langages non connus de listings
\lstdefinelanguage{yaml}{
  keywords={true,false,null,on,off,needs,runs-on,uses,with,env,jobs,steps,name,
            strategy,matrix,inputs,secrets,if,run,outputs,permissions},
  sensitive=false,
  comment=[l]{\#},
  morestring=[b]",
  morestring=[b]'
}
\lstdefinelanguage{docker}{
  keywords={FROM,RUN,COPY,ADD,ARG,ENV,WORKDIR,USER,ENTRYPOINT,CMD,LABEL,SHELL,AS},
  sensitive=true,
  comment=[l]{\#},
  morestring=[b]"
}
\lstdefinelanguage{hcl}{
  keywords={resource,module,variable,output,provider,locals,data,terraform,for_each,source},
  sensitive=true,
  comment=[l]{\#},
  morestring=[b]"
}

% ---------------------------------------------------------------------------
%  Bibliographie
% ---------------------------------------------------------------------------
\usepackage[backend=biber,style=numeric,sorting=nyt,maxbibnames=99,giveninits=true]{biblatex}
\addbibresource{biblio.bib}

% ---------------------------------------------------------------------------
%  Acronymes / glossaire
% ---------------------------------------------------------------------------
\usepackage[acronym,nonumberlist,nopostdot,nomain]{glossaries}
\makeglossaries

% ---------------------------------------------------------------------------
%  Notes de rédaction  (mettre \setboolean{draftmode}{true} pour la version finale)
% ---------------------------------------------------------------------------
\usepackage{ifthen}
\newboolean{draftmode}
\setboolean{draftmode}{true}

\usepackage[colorinlistoftodos,prependcaption,textsize=footnotesize,french]{todonotes}
\ifthenelse{\boolean{draftmode}}{}{\presetkeys{todonotes}{disable}{}}

% Raccourcis de rédaction
\newcommand{\aecrire}[1]{\todo[inline,color=orange!25]{À RÉDIGER : #1}}
\newcommand{\chiffre}[1]{\todo[inline,color=red!25]{CHIFFRE À MESURER : #1}\textbf{40 Milliards d'euros}}
\newcommand{\verifier}[1]{\todo[color=blue!25]{VÉRIFIER : #1}}

% ---------------------------------------------------------------------------
%  ANONYMISATION
%  Tant que l'autorisation de publication n'est pas obtenue, basculer
%  \setboolean{confidentiel}{true} : tous les noms sensibles sont remplacés.
% ---------------------------------------------------------------------------
\newboolean{confidentiel}
\setboolean{confidentiel}{false}

\ifthenelse{\boolean{confidentiel}}{%
  \newcommand{\Entreprise}{l'entreprise}
  \newcommand{\Produit}{le produit P}
  \newcommand{\ProduitCourt}{P}
  \newcommand{\Firmware}{le firmware applicatif}
  \newcommand{\Carte}{la carte cible}
  \newcommand{\SrvCoverity}{\texttt{<serveur-sast-interne>}}
  \newcommand{\SrvBDBA}{\texttt{<serveur-sca-interne>}}
  \newcommand{\SrvArtifacts}{\texttt{<artifactory-interne>}}
  \newcommand{\SrvGit}{\texttt{<github-entreprise>}}
}{%
  \newcommand{\Entreprise}{Schneider Electric}
  \newcommand{\Produit}{le PM77xx}
  \newcommand{\ProduitCourt}{PM77xx}
  \newcommand{\Firmware}{SDP}
  \newcommand{\Carte}{bru99379-kp0}
  \newcommand{\SrvCoverity}{\texttt{coverity-embu.se.com}}
  \newcommand{\SrvBDBA}{\texttt{bdba.schneider-electric.com}}
  \newcommand{\SrvArtifacts}{\texttt{global-artifacts.se.com}}
  \newcommand{\SrvGit}{\texttt{github.schneider-electric.com}}
}

% ---------------------------------------------------------------------------
%  Hyperliens (à charger en dernier ou presque)
% ---------------------------------------------------------------------------
\usepackage[hidelinks,breaklinks=true]{hyperref}
\hypersetup{
  pdftitle={Conception et mise en oeuvre d'une infrastructure CI/CD pour le firmware d'un produit de comptage d'energie},
  pdfsubject={Memoire de Master Ingenierie des Systemes et Reseaux},
  pdfkeywords={CI/CD, DevOps, systemes embarques, Yocto, GitHub Actions, analyse statique, SBOM, Infrastructure as Code}
}
\usepackage[french]{cleveref}
